\section{一次循环的工作原理}

\subsection{菜谱类比}

作者用一个直观的类比来解释 autoresearch 的单次循环：

\begin{knowledgebox}{菜谱类比——理解 AutoResearch 循环}
你有一个菜谱，十次里有七次做得不错，另外三次总是不对劲——酱汁太淡，或者调味不对。

你不是从头重写整个菜谱，而是：
\begin{enumerate}
    \item 换一种配料
    \item 用这个改动做十次
    \item 变好了？留下这个改动
    \item 变差了？换回原来的配料
    \item 改下一个东西，再做十次
\end{enumerate}

经过 50 轮这样的迭代，你的菜谱十次里有 9.5 次都能成功。
\end{knowledgebox}

\subsection{映射到 Skill 优化}

这个类比精确对应到 skill 优化的每个环节：

\begin{itemize}
    \item \textbf{菜谱（recipe）} = 你的 skill prompt
    \item \textbf{做菜（cooking）} = 运行这个 skill
    \item \textbf{试味（tasting）} = 给输出结果打分
\end{itemize}

而你唯一需要提供的，就是\textbf{打分标准}（scoring criteria）。

\subsection{本章小结}

单次循环的逻辑是：改一个点 $\to$ 测试多次 $\to$ 根据分数决定保留或撤销。通过大量循环的累积，skill 的整体质量从"七成靠谱"逐步提升到"九成五靠谱"。这是一个经典的贪心搜索过程。

